\documentclass[11pt]{article}
\usepackage[margin=1in]{geometry}
\usepackage{amsmath,amssymb,amsthm,hyperref}
\newtheorem{theorem}{Theorem}
\newtheorem{lemma}{Lemma}
\newtheorem{corollary}{Corollary}
\newcommand{\tr}{\operatorname{tr}}
\newcommand{\Rea}{\operatorname{Re}}
\newcommand{\Ima}{\operatorname{Im}}
\newcommand{\HH}{\mathcal H}
\title{Trace-small imaginary parts preserve spectral symbols}
\author{Principia Math in Progress}
\date{September 12, 2026}
\begin{document}
\maketitle

\section*{Task 1. Solve the real-part spectral-symbol conjecture}
We prove Conjecture~2 on page~6 of Barbarino's
\emph{Conjectures on Perturbations of Hermitian Sequences}
\cite{Barbarino}, corresponding to MathDB~340128.
For a matrix \(X\), put
\(\Rea X=(X+X^*)/2\) and \(\Ima X=(X-X^*)/(2i)\);
\(\|X\|_1\) denotes the sum of its singular values.
For \(n\times n\) matrices, write \(X_n\sim_\lambda f\) when
\[
 \frac1n\sum_{j=1}^n F(\lambda_j(X_n))
 \longrightarrow \frac1{|D|}\int_D F(f(x))\,dx
 \qquad(F\in C_c(\mathbb C)),
 \tag{1}
\]
where \(f:D\to\mathbb C\) is measurable and \(0<|D|<\infty\).
Eigenvalues are counted with algebraic multiplicity.

\begin{theorem}\label{main}
If \(X_n\sim_\lambda f\) and \(\|\Ima X_n\|_1=o(n)\), then
\(\Rea X_n\sim_\lambda f\). In particular \(f\) is real almost everywhere.
No bound on \(\|X_n\|\) is required.
\end{theorem}

The main estimate is a finite-dimensional weak trace-norm inequality
for triangular truncation. Such inequalities are classical; see, for
example, the estimate cited in the proof of Lemma~3.2 in \cite{CPSZ}.
We include an elementary determinant proof, so the argument does not
depend on that result.

For Hermitian \(B=(b_{jk})\), let \(U(B)\) be its strictly upper triangular
part and define the Hermitian matrix
\[
 \HH(B)=i\bigl(U(B)-U(B)^*\bigr).
\]
For any matrix \(A\), let \(N_t(A)=\#\{j:s_j(A)>t\}\), for \(t>0\).

\begin{lemma}\label{weak}
For every Hermitian matrix \(B\) of any size,
\[
 N_t(\HH(B))\le \frac{4}{\log 2}\frac{\|B\|_1}{t}.
 \tag{2}
\]
\end{lemma}
\begin{proof}
First suppose \(B\ge0\) and write \(H=\HH(B)\).
The matrix \(B+iH=\operatorname{diag}(B)+2U(B)^*\) is lower triangular.
Hence
\[
 \left|\det(I+B/t+iH/t)\right|
 =\prod_j(1+b_{jj}/t)
 \le \exp(\tr B/t).
 \tag{3}
\]
We also have
\[
 |\det(I+B/t+iH/t)|\ge |\det(I+iH/t)|.
 \tag{4}
\]
Indeed, put \(P=B/t\ge0\), \(Q=H/t=Q^*\), and
\(Z=P^{1/2}(I+iQ)^{-1}P^{1/2}\).
Since
\[
 \Rea Z=P^{1/2}(I+Q^2)^{-1}P^{1/2}\ge0,
\]
every eigenvalue \(z\) of \(Z\) satisfies \(\Rea z\ge0\).
Thus \(|\det(I+Z)|=\prod_z|1+z|\ge1\).
Factoring \(I+P+iQ\) and using
\(\det(I+AB)=\det(I+BA)\) proves (4).

Combining (3)--(4) and taking logarithms gives
\[
 \frac12\sum_j\log\bigl(1+\lambda_j(H)^2/t^2\bigr)
 \le \frac{\tr B}{t}.
\]
Each term with \(|\lambda_j(H)|>t\) exceeds \((\log 2)/2\).
Consequently \(N_t(H)\le 2\tr B/(t\log 2)\).

For general \(B\), use its positive and negative spectral parts:
\(B=B_+-B_-\), with \(B_\pm\ge0\) and
\(\tr B_++\tr B_-=\|B\|_1\). Linearity and the singular-value
counting inequality give
\[
 N_t(\HH(B))
 \le N_{t/2}(\HH(B_+))+N_{t/2}(\HH(B_-))
 \le \frac{4\|B\|_1}{t\log2}.
\]
For completeness, the counting inequality follows by splitting each
summand into its singular directions above \(t/2\) and a remainder
of norm at most \(t/2\). The sum of the former pieces has rank at most
the sum of the two counts, and the total remainder has norm at most \(t\).
\end{proof}

\begin{lemma}\label{stability}
Suppose \(A_n,D_n\) are Hermitian, of size \(n\), and
\(N_t(A_n-D_n)=o(n)\) for every fixed \(t>0\).
Then, for every \(g\in C_c(\mathbb R)\),
\[
 \frac1n\tr g(A_n)-\frac1n\tr g(D_n)\longrightarrow0.
 \tag{5}
\]
\end{lemma}
\begin{proof}
Fix \(t>0\). Spectrally truncate \(A_n-D_n=R_n+E_n\), retaining in
\(R_n\) only the eigenvalues whose absolute values exceed \(t\).
Then both pieces are Hermitian, \(\operatorname{rank}R_n=o(n)\),
and \(\|E_n\|\le t\).
The min--max principle bounds the difference of the empirical
distribution functions of \(D_n+R_n\) and \(D_n\) by
\(\operatorname{rank}R_n/n\).
For \(g\in C_c^1(\mathbb R)\), integration by parts therefore bounds
the difference of their normalized traces by
\[
 \frac{\operatorname{rank}R_n}{n}\int_{\mathbb R}|g'(x)|\,dx.
\]
The ordered eigenvalues of \(A_n=D_n+R_n+E_n\) and \(D_n+R_n\)
differ by at most \(t\), again by min--max. Their normalized
\(g\)-traces differ by at most
\(\omega_g(t)=\sup_{|x-y|\le t}|g(x)-g(y)|\).
First let \(n\to\infty\), then \(t\downarrow0\).
Uniform approximation by \(C_c^1(\mathbb R)\) extends the conclusion
to all \(g\in C_c(\mathbb R)\).
\end{proof}

\begin{proof}[Proof of Theorem~\ref{main}]
Take a unitary Schur decomposition
\[
 V_n^*X_nV_n=L_n+T_n,\qquad
 L_n=\operatorname{diag}(\lambda_1(X_n),\ldots,\lambda_n(X_n)),
\]
where \(T_n\) is strictly upper triangular.
Put \(K_n=\Ima(L_n+T_n)\) and
\(D_n=\operatorname{diag}(\Rea\lambda_j(X_n))\).
Then \(\|K_n\|_1=\|\Ima X_n\|_1=o(n)\), and a direct entry calculation gives
\[
 V_n^*(\Rea X_n)V_n=D_n+\HH(K_n).
 \tag{6}
\]
Indeed \(U(K_n)=T_n/(2i)\), so
\(\HH(K_n)=(T_n+T_n^*)/2\).
Lemma~\ref{weak} implies
\(N_t(\HH(K_n))=o(n)\) for every \(t>0\).
Lemma~\ref{stability} then shows that \(\Rea X_n\) and \(D_n\)
have asymptotically equal normalized traces against every continuous
compactly supported function on the real line.

The diagonal entries of \(K_n\) are \(\Ima\lambda_j(X_n)\).
The diagonal trace-norm bound gives
\[
 \sum_j|\Ima\lambda_j(X_n)|\le\|K_n\|_1=o(n).
 \tag{7}
\]
To verify this bound directly, choose a diagonal matrix \(S\) whose
entries are the signs of these real diagonal entries. Then
\(\sum_j|(K_n)_{jj}|=\tr(SK_n)\le\|K_n\|_1\), since \(\|S\|\le1\).

Given \(F\in C_c(\mathbb C)\), it is uniformly continuous on \(\mathbb C\).
For every \(\delta>0\), (7) bounds the difference between its eigenvalue
averages on \(X_n\) and \(D_n\) by
\[
 \omega_F(\delta)
 +2\|F\|_\infty\frac{\|K_n\|_1}{n\delta}.
 \tag{8}
\]
Let \(n\to\infty\) and then \(\delta\downarrow0\).
Thus \(D_n\sim_\lambda f\); applying (5) to \(F|_{\mathbb R}\)
and using (6) proves the required conclusion.
Finally the limiting probability measure in (1) is supported on
\(\mathbb R\): every nonnegative compactly supported continuous test
function away from \(\mathbb R\) has integral zero.
An exhaustion of \(\mathbb C\setminus\mathbb R\) by compact subsets
shows that \(f\) is real almost everywhere.
\end{proof}

\begin{corollary}\label{forward}
If \(A_n=A_n^*\sim_\lambda f\) and \(\|Y_n\|_1=o(n)\), then
\(A_n+Y_n\sim_\lambda f\).
\end{corollary}
\begin{proof}
Set \(X_n=A_n+Y_n\). Since
\(\|\Rea Y_n\|_1,\|\Ima Y_n\|_1\le\|Y_n\|_1\),
Lemma~\ref{stability} gives \(\Rea X_n\sim_\lambda f\).
The Schur argument (6)--(8) compares the eigenvalue averages of
\(X_n\) and \(\Rea X_n\) without assuming that either already has
a symbol. Their difference tends to zero, proving the claim.
\end{proof}

This also proves the companion perturbation conjecture in
\cite{Barbarino,BS}. It is recorded here as a corollary of the same
argument and receives no separate campaign credit.

\section*{Task 2. Determine the sharp triangular counting constant}
Find the least dimension-independent \(C\) for which
\(N_t(\HH(B))\le C\|B\|_1/t\) for all Hermitian \(B\).
Describe families approaching the optimal constant.

\section*{Task 3. Find optimal rates for spectral convergence}
Under a uniform density bound for the limiting real eigenvalue law,
determine the best rate in Kolmogorov or L\'evy distance in terms of
\(\|\Ima X_n\|_1/n\) and the initial spectral approximation error.

\section*{Task 4. Classify other spectral supports with this stability}
Characterize closed sets \(E\subset\mathbb C\) for which every normal
sequence with spectra in \(E\) keeps its spectral law under arbitrary
trace-norm \(o(n)\) perturbations. Affine real lines satisfy this
property; the circular counterexample in \cite{Barbarino} shows that
it fails for some other supports.

\section*{Task 5. Find weaker sufficient hypotheses on imaginary parts}
Identify natural conditions on the singular-value tails of
\(\Ima X_n\), weaker than trace-norm \(o(n)\), that still guarantee
equality of the spectral laws of \(X_n\) and \(\Rea X_n\).
Require conditions that also control the triangular transform, since
zero singular-value distribution alone is insufficient.

\begin{thebibliography}{9}
\bibitem{Barbarino}
G. Barbarino, \emph{Conjectures on Perturbations of Hermitian Sequences},
arXiv:1808.05555v1, Conjectures~1--2.
\url{https://arxiv.org/abs/1808.05555}.
\bibitem{CPSZ}
M. Caspers, D. Potapov, F. Sukochev and D. Zanin,
\emph{Weak type estimates for the absolute value mapping},
arXiv:1309.3378; Journal of Operator Theory (2015).
\url{https://arxiv.org/abs/1309.3378}.
\bibitem{BS}
G. Barbarino and S. Serra-Capizzano,
\emph{Non-Hermitian perturbations of Hermitian matrix-sequences and
applications to the spectral analysis of the numerical approximation
of partial differential equations}, Numerical Linear Algebra with
Applications (2020), Conjecture~1.
\url{https://giovannibarbarino.github.io/doc/articles/NHperturbation.pdf}.
\end{thebibliography}
\end{document}
